% !TEX root = ../main.tex
%% Chapitre 13 — Orthogonalité et distances dans l'espace
%% Sources : frontend/exercices/terminal_spé_Maths/exercice_produit_scalaire_terminal.txt
%%           frontend/cours/resume/terminal/resume_representations_parametriques_terminal.txt (équations cartésiennes)
%% Tous les calculs ont été vérifiés par programme (calcul exact).
%% Exercices 1, 6, 7, 8 et 12 : exercices inédits rédigés pour le livre (numérotation de l'édition 2027-2028).
%% Édition 2027-2028 : théorème « Plans perpendiculaires » et exercice 9 nouveaux (programme de la rentrée 2027).

\begin{ouverture}[orthogonalite]{13}{Orthogonalité et distances}{Partie 2 · Géométrie dans l'espace}
  {Produit scalaire\enspace·\enspace Orthogonalité\enspace·\enspace Vecteur normal\enspace·\enspace Équation cartésienne\enspace·\enspace
   Projeté orthogonal}
  \objectif{Calculer un produit scalaire dans l'espace, avec ou sans coordonnées, et en déduire une longueur ou un angle.}
  \objectif{Démontrer l'orthogonalité de deux droites, d'une droite et d'un plan.}
  \objectif{Déterminer un vecteur normal et une équation cartésienne d'un plan.}
  \objectif{Déterminer le projeté orthogonal d'un point sur un plan ou sur une droite, et la distance correspondante.}
  \prerequis{Produit scalaire dans le plan (Première).}
  \prerequis{Vecteurs, bases et repères de l'espace (chapitres 11 et 12).}
\end{ouverture}
\miseEnPageSouple

%% =====================================================================
%%  L'ESSENTIEL
%% =====================================================================
\essentiel

\essSection{1}{Produit scalaire dans l'espace}

\noindent
\begin{minipage}[c]{0.58\linewidth}
\begin{definition}[Produit scalaire]
\raggedright
Soit $\vec u=\vect{AB}$ et $\vec v=\vect{AC}$. Les points $A$, $B$, $C$ sont dans un même plan, et l'on définit $\vec u\cdot\vec v$
comme dans ce plan :
\[\vec u\cdot\vec v=\norme{\vec u}\times\norme{\vec v}\times\cos\widehat{BAC}=\tfrac12\left(\norme{\vec u+\vec v}^2-\norme{\vec u}^2-\norme{\vec v}^2\right).\]
Si $H$ est le projeté orthogonal de $C$ sur $(AB)$, alors $\vect{AB}\cdot\vect{AC}=\vect{AB}\cdot\vect{AH}$.
\end{definition}
\end{minipage}\hfill
\begin{minipage}[c]{0.38\linewidth}
\centering
\begin{tikzpicture}[scale=1]
  \fill[otNavyTint] (-0.5,-0.55) -- (3.6,-0.55) -- (4.5,1.9) -- (0.4,1.9) -- cycle;
  \coordinate (A) at (0.2,0); \coordinate (B) at (3.2,0.25); \coordinate (C) at (2.1,1.55);
  \coordinate (H) at ($(A)!(C)!(B)$);
  \draw[otCache] (C) -- (H);
  \draw[otVecN] (A) -- (B); \draw[otVec] (A) -- (C);
  \coordinate (h1) at ($(H)!0.18cm!(A)$); \coordinate (h2) at ($(H)!0.18cm!(C)$);
  \draw[otMuted,line width=0.5pt] (h1) -- ($(h1)+(h2)-(H)$) -- (h2);
  \foreach \p/\a in {A/north east,B/north west,C/south,H/north}{\node[otSommet,anchor=\a] at (\p) {$\p$};}
\end{tikzpicture}\\[0.5mm]
\legende{$\vect{AB}\cdot\vect{AC}=AB\times AH$ ici.}
\end{minipage}

\begin{propriete}[Règles de calcul et repère orthonormé]
Le produit scalaire est symétrique et bilinéaire : $\vec u\cdot\vec v=\vec v\cdot\vec u$,
$\vec u\cdot(\vec v+\vec w)=\vec u\cdot\vec v+\vec u\cdot\vec w$, $(k\vec u)\cdot\vec v=k(\vec u\cdot\vec v)$, et
$\vec u\cdot\vec u=\norme{\vec u}^2$. Dans un repère \textbf{orthonormé}, si $\vec u\,(x\pv y\pv z)$ et $\vec v\,(x'\pv y'\pv z')$ :
\[\vec u\cdot\vec v=xx'+yy'+zz',\qquad \norme{\vec u}=\sqrt{x^2+y^2+z^2},\qquad
AB=\sqrt{(x_B-x_A)^2+(y_B-y_A)^2+(z_B-z_A)^2}.\]
\end{propriete}

\essSection{2}{Orthogonalité}

\noindent
\begin{minipage}[c]{0.6\linewidth}
\begin{propriete}[Vecteurs, droites et plans orthogonaux]
\raggedright
\begin{itemize}[leftmargin=1.2em,itemsep=0.4mm,topsep=0.3mm]
  \item $\vec u$ et $\vec v$ sont \textbf{orthogonaux} si et seulement si $\vec u\cdot\vec v=0$.
  \item Deux droites sont orthogonales si leurs vecteurs directeurs le sont ; elles ne sont pas forcément sécantes. Orthogonales et
    sécantes, elles sont \textbf{perpendiculaires}.
  \item Une droite est \textbf{orthogonale à un plan} si et seulement si elle est orthogonale à deux droites sécantes de ce plan.
    Elle est alors orthogonale à toutes les droites du plan.
\end{itemize}
\end{propriete}
\end{minipage}\hfill
\begin{minipage}[c]{0.36\linewidth}
\centering
\begin{tikzpicture}[otEspace,scale=2.3]
  \otCube
  \draw[otGreen,line width=2pt] (A) -- (E);
  \draw[otOrange,line width=2pt] (H) -- (G);
\end{tikzpicture}\\[0.5mm]
\legende{$(AE)$ et $(HG)$ sont orthogonales, mais non sécantes. $(AE)$ est orthogonale au plan $(ABC)$.}
\end{minipage}

\essSection{3}{Vecteur normal et équation cartésienne d'un plan}

\noindent
\begin{minipage}[c]{0.58\linewidth}
\begin{theoreme}[Équation cartésienne]
\raggedright
Un vecteur $\vec n\neq\vec0$ est \textbf{normal} à un plan $\mathcal{P}$ s'il est orthogonal à deux vecteurs non colinéaires de
$\mathcal{P}$. Dans un repère orthonormé, si $\vec n\,(a\pv b\pv c)$ est normal à $\mathcal{P}$ et $A\in\mathcal{P}$ :
\[M\in\mathcal{P}\iff\vect{AM}\cdot\vec n=0\iff ax+by+cz+d=0.\]
Réciproquement, si $(a\pv b\pv c)\neq(0\pv0\pv0)$, l'ensemble des points vérifiant $ax+by+cz+d=0$ est un plan de vecteur normal
$\vec n\,(a\pv b\pv c)$.
\end{theoreme}
\end{minipage}\hfill
\begin{minipage}[c]{0.38\linewidth}
\centering
\begin{tikzpicture}[scale=1]
  \fill[otNavyTint] (-0.6,-0.6) -- (3.3,-0.6) -- (4.2,1.4) -- (0.3,1.4) -- cycle;
  \draw[otNavy!45,line width=0.5pt] (-0.6,-0.6) -- (3.3,-0.6) -- (4.2,1.4) -- (0.3,1.4) -- cycle;
  \coordinate (A) at (1.2,0.3); \coordinate (M) at (2.9,0.9);
  \draw[otVec] (A) -- (1.2,2.3) node[left,font=\small] {$\vec n$};
  \draw[otVecN] (A) -- (M);
  \draw[otMuted,line width=0.5pt] (1.2,0.5) -- (1.389,0.567) -- (1.389,0.367);
  \node[otPt] at (A) {}; \node[otSommet,anchor=north east] at (A) {$A$};
  \node[otPt] at (M) {}; \node[otSommet,anchor=north west] at (M) {$M$};
  \node[font=\small,text=otNavy] at (3.6,-0.3) {$\mathcal{P}$};
\end{tikzpicture}\\[0.5mm]
\legende{$M\in\mathcal{P}$ si et seulement si $\vect{AM}\perp\vec n$.}
\end{minipage}

\begin{demonstration}[De l'orthogonalité à l'équation]
Avec $A(x_A\pv y_A\pv z_A)$ et $M(x\pv y\pv z)$ : $\vect{AM}\cdot\vec n=a(x-x_A)+b(y-y_A)+c(z-z_A)$. Cette quantité est nulle si et
seulement si $ax+by+cz+d=0$, en posant $d=-(ax_A+by_A+cz_A)$.
\end{demonstration}

\begin{savoirfaire}[Déterminer une équation cartésienne d'un plan]
\textbf{Exemple.} Plan passant par $A(1\pv0\pv2)$ et de vecteur normal $\vec n\,(2\pv-1\pv1)$ : une équation est
$2x-y+z+d=0$, et $A$ vérifie $2-0+2+d=0$, d'où $d=-4$. Le plan a pour équation $\boxed{2x-y+z-4=0}$. Pour un plan défini par trois
points, on cherche d'abord un vecteur $\vec n\,(a\pv b\pv c)$ orthogonal à $\vect{AB}$ et à $\vect{AC}$ (système de deux équations).
\end{savoirfaire}

\begin{propriete}[Positions relatives de deux plans]
Deux plans de vecteurs normaux $\vec{n_1}$ et $\vec{n_2}$ sont \textbf{parallèles} si et seulement si $\vec{n_1}$ et $\vec{n_2}$ sont
colinéaires. Une droite de vecteur directeur $\vec u$ est orthogonale à un plan si et seulement si $\vec u$ est colinéaire à un
vecteur normal du plan.
\end{propriete}

\noindent
\begin{minipage}[c]{0.6\linewidth}
\begin{theoreme}[Plans perpendiculaires\enspace\otNouveau]
\raggedright
Deux plans $\mathcal{P}_1$ et $\mathcal{P}_2$ sont \textbf{perpendiculaires} si l'un contient une droite orthogonale à l'autre.
Si $\vec{n_1}$ et $\vec{n_2}$ sont des vecteurs normaux à $\mathcal{P}_1$ et à $\mathcal{P}_2$ :
\[\mathcal{P}_1\perp\mathcal{P}_2\iff\vec{n_1}\cdot\vec{n_2}=0.\]
\end{theoreme}
\end{minipage}\hfill
\begin{minipage}[c]{0.36\linewidth}
\centering
\begin{tikzpicture}[scale=1]
  \fill[otNavyTint] (-0.4,-0.6) -- (3.2,-0.6) -- (4.1,0.6) -- (0.5,0.6) -- cycle;
  \draw[otNavy!45,line width=0.5pt] (-0.4,-0.6) -- (3.2,-0.6) -- (4.1,0.6) -- (0.5,0.6) -- cycle;
  \fill[otGreen,opacity=0.22] (1.2,-0.4) -- (1.8,0.4) -- (1.8,2.4) -- (1.2,1.6) -- cycle;
  \draw[otGreenDark,line width=0.6pt] (1.2,-0.4) -- (1.8,0.4) -- (1.8,2.4) -- (1.2,1.6) -- cycle;
  \draw[otVecN] (1.5,0) -- (1.5,1.15) node[left,font=\small,text=otNavy,xshift=-1pt] {$\vec{n_1}$};
  \draw[otVec] (1.5,1.4) -- (2.65,1.4) node[right,font=\small,text=otGreenDark] {$\vec{n_2}$};
  \node[font=\small,text=otNavy] at (3.45,-0.3) {$\mathcal{P}_1$};
  \node[font=\small,text=otGreenDark] at (1.05,2.05) {$\mathcal{P}_2$};
\end{tikzpicture}\\[0.5mm]
\legende{$\mathcal{P}_2$ contient la droite dirigée par $\vec{n_1}$.}
\end{minipage}

\begin{demonstration}
Si $\vec{n_1}\cdot\vec{n_2}=0$, le vecteur $\vec{n_1}$ est orthogonal à $\vec{n_2}$ : la droite passant par un point de
$\mathcal{P}_2$ et dirigée par $\vec{n_1}$ est contenue dans $\mathcal{P}_2$, et elle est orthogonale à $\mathcal{P}_1$.
Réciproquement, si $\mathcal{P}_2$ contient une droite orthogonale à $\mathcal{P}_1$, un vecteur directeur de cette droite est
colinéaire à $\vec{n_1}$ et orthogonal à $\vec{n_2}$, donc $\vec{n_1}\cdot\vec{n_2}=0$.
\end{demonstration}

\essSection{4}{Projeté orthogonal et distance}

\noindent
\begin{minipage}[c]{0.58\linewidth}
\begin{definition}[Projeté orthogonal]
\raggedright
Le projeté orthogonal d'un point $M$ sur un plan $\mathcal{P}$ est le point $H$ de $\mathcal{P}$ tel que $(MH)$ soit orthogonale à
$\mathcal{P}$. Pour tout point $Q$ de $\mathcal{P}$, $MQ^2=MH^2+HQ^2\geq MH^2$ : $MH$ est la plus petite distance de $M$ à un point
du plan. C'est la \textbf{distance} de $M$ à $\mathcal{P}$.\\[1mm]
Le projeté orthogonal de $M$ sur une droite $d$ de vecteur directeur $\vec u$ est le point $H$ de $d$ tel que $\vect{MH}\cdot\vec u=0$.
\end{definition}
\end{minipage}\hfill
\begin{minipage}[c]{0.38\linewidth}
\centering
\begin{tikzpicture}[scale=1]
  \fill[otNavyTint] (-0.6,-0.6) -- (3.3,-0.6) -- (4.2,1.4) -- (0.3,1.4) -- cycle;
  \draw[otNavy!45,line width=0.5pt] (-0.6,-0.6) -- (3.3,-0.6) -- (4.2,1.4) -- (0.3,1.4) -- cycle;
  \coordinate (H) at (1.3,0.4); \coordinate (M) at (1.3,2.3); \coordinate (Q) at (3.1,0.1);
  \draw[otGreenDark,line width=1pt] (M) -- (H);
  \draw[otMuted,line width=0.7pt] (M) -- (Q) -- (H);
  \draw[otMuted,line width=0.5pt] (1.3,0.6) -- (1.48,0.57) -- (1.48,0.37);
  \node[otPtV] at (M) {}; \node[otSommet,anchor=south] at (M) {$M$};
  \node[otPt] at (H) {}; \node[otSommet,anchor=north east] at (H) {$H$};
  \node[otPt] at (Q) {}; \node[otSommet,anchor=north] at (Q) {$Q$};
  \node[font=\small,text=otNavy] at (3.7,1.05) {$\mathcal{P}$};
\end{tikzpicture}\\[0.5mm]
\legende{Le triangle $MHQ$ est rectangle en $H$.}
\end{minipage}

\begin{savoirfaire}[Calculer la distance d'un point à un plan]
On paramètre la droite passant par $M$ et dirigée par $\vec n$ : ses points s'écrivent $M+t\vec n$. On reporte dans l'équation du plan
pour trouver $t$, donc $H$ ; la distance cherchée est $MH=|t|\times\norme{\vec n}$.
\end{savoirfaire}

\begin{exemple}[Distance de l'origine à un plan]
Soit $\mathcal{P}:x+2y+2z-9=0$, de vecteur normal $\vec n\,(1\pv2\pv2)$. Les points de la droite passant par $O$ et dirigée par $\vec n$
s'écrivent $(t\pv2t\pv2t)$. En reportant : $t+4t+4t-9=0$, soit $t=1$. Le projeté de $O$ est $H(1\pv2\pv2)$, et la distance de $O$ à
$\mathcal{P}$ vaut $OH=\sqrt{1+4+4}=3$.
\end{exemple}

\begin{attention}[Pièges fréquents]
\begin{itemize}[label={\color{otOrange}\textbullet},leftmargin=1.4em,itemsep=0.6mm,topsep=0.5mm]
  \item La formule $xx'+yy'+zz'$ n'est valable que dans un repère \textbf{orthonormé}.
  \item Deux droites orthogonales ne sont pas forcément sécantes ; « perpendiculaires » signifie orthogonales \emph{et} sécantes.
  \item Un plan admet une infinité de vecteurs normaux et d'équations cartésiennes, toutes proportionnelles.
  \item Pour qu'une droite soit orthogonale à un plan, il ne suffit pas qu'elle soit orthogonale à une seule droite du plan.
\end{itemize}
\end{attention}

%% =====================================================================
%%  EXERCICES
%% =====================================================================
\exercices
\rubrique{Découverte}

\begin{exercice}[1]{Produit scalaire en repère orthonormé}
L'espace est muni d'un repère orthonormé. On considère $\vec u\,(2\pv-1\pv3)$, $\vec v\,(1\pv4\pv2)$ et $\vec w\,(-3\pv0\pv2)$.
\begin{questions}
  \item Calculer $\vec u\cdot\vec v$ et $\vec u\cdot\vec w$. Que peut-on dire de $\vec u$ et $\vec w$ ?
  \item Calculer $\norme{\vec u}$ et $\norme{\vec v}$, puis une valeur approchée au degré près de l'angle formé par $\vec u$ et $\vec v$.
  \item Vérifier sur cet exemple l'égalité $\norme{\vec u+\vec v}^2=\norme{\vec u}^2+2\,\vec u\cdot\vec v+\norme{\vec v}^2$.
  \item Calculer la distance $AB$, où $A(1\pv0\pv2)$ et $B(3\pv1\pv0)$.
\end{questions}
\end{exercice}

\begin{exercice}[1]{Produit scalaire dans un cube}
\begin{codecote}[0.6]
On considère un cube $ABCDEFGH$ d'arête $1$.
\begin{questions}
  \item Sans repère, en décomposant $\vect{AG}=\vect{AC}+\vect{CG}$, calculer $\vect{AC}\cdot\vect{AG}$.
  \item Retrouver ce résultat dans le repère orthonormé $\big(A\pv\vect{AB},\vect{AD},\vect{AE}\big)$.
  \item Calculer de même $\vect{AB}\cdot\vect{AG}$ et $\vect{EC}\cdot\vect{BH}$.
  \item En déduire une valeur approchée de l'angle $\widehat{CAG}$.
\end{questions}
\cote
\centering
\begin{tikzpicture}[otEspace,scale=2.3]
  \otCube
  \draw[otGreenDark,line width=0.9pt,dash pattern=on 2.2pt off 1.6pt] (A) -- (C) (A) -- (G);
  \draw[otGreenDark,line width=0.9pt] (C) -- (G);
\end{tikzpicture}
\end{codecote}
\end{exercice}

\begin{exercice}[1]{Vecteur normal et équation d'un plan}
L'espace est muni d'un repère orthonormé. On considère les points $A(2\pv1\pv-1)$, $B(0\pv3\pv2)$, $C(-1\pv0\pv1)$ et le vecteur
$\vec n\,(7\pv-5\pv8)$.
\begin{questions}
  \item Justifier que les points $A$, $B$ et $C$ définissent un plan.
  \item Montrer que $\vec n$ est un vecteur normal au plan $(ABC)$.
  \item En déduire une équation cartésienne du plan $(ABC)$. Vérifier que les coordonnées de $B$ et de $C$ la satisfont.
\end{questions}
\end{exercice}

\rubrique{Entraînement}

\begin{exercice}[2]{Produits scalaires dans un tétraèdre régulier}
\begin{codecote}[0.6]
$ABCD$ est un tétraèdre régulier d'arête $a$ : ses six arêtes mesurent $a$ et ses faces sont des triangles équilatéraux. $I$, $J$ et $K$
sont les milieux de $[AB]$, $[AC]$ et $[AD]$.
\begin{questions}
  \item Calculer $\vect{AB}\cdot\vect{AC}$.
  \item Calculer $\vect{BI}\cdot\vect{AJ}$, puis $\vect{IK}\cdot\vect{JK}$.
  \item Montrer que $\vect{BC}\cdot\vect{AD}=0$ en écrivant $\vect{BC}=\vect{AC}-\vect{AB}$. En déduire $\vect{IJ}\cdot\vect{AD}$.
  \item Que peut-on dire de deux arêtes opposées d'un tétraèdre régulier ?
\end{questions}
\cote
\centering
\begin{tikzpicture}[scale=1.3]
  \coordinate (A) at (0,0); \coordinate (B) at (1.55,-0.85); \coordinate (C) at (2.95,0.1); \coordinate (D) at (1.45,2.2);
  \coordinate (I) at ($(A)!0.5!(B)$); \coordinate (J) at ($(A)!0.5!(C)$); \coordinate (K) at ($(A)!0.5!(D)$);
  \draw[otVu] (A) -- (B) -- (C) -- (D) -- cycle; \draw[otVu] (B) -- (D);
  \draw[otCache] (A) -- (C);
  \draw[otGreenDark,line width=0.9pt] (I) -- (K);
  \draw[otGreenDark,line width=0.9pt,dash pattern=on 2.2pt off 1.6pt] (I) -- (J) -- (K);
  \foreach \p/\a in {I/north east,J/north west,K/south east}{\node[otPtV] at (\p) {}; \node[otSommet,anchor=\a] at (\p) {$\p$};}
  \node[otSommet,anchor=east] at (A) {$A$}; \node[otSommet,anchor=north] at (B) {$B$};
  \node[otSommet,anchor=west] at (C) {$C$}; \node[otSommet,anchor=south] at (D) {$D$};
\end{tikzpicture}
\end{codecote}
\end{exercice}

\begin{exercice}[2]{Calculer un angle dans un cube}
\begin{codecote}[0.6]
$ABCDEFGH$ est un cube d'arête $1$. On note $M$ le milieu de $[AD]$ et $N$ le milieu de $[EF]$. On cherche une mesure de l'angle
$\widehat{HMN}$.
\begin{questions}
  \item Donner les coordonnées de $H$, $M$ et $N$ dans le repère orthonormé $\big(A\pv\vect{AB},\vect{AD},\vect{AE}\big)$.
  \item Calculer $\vect{MH}\cdot\vect{MN}$, puis les longueurs $MH$ et $MN$.
  \item En déduire $\cos\widehat{HMN}$, puis une mesure de l'angle $\widehat{HMN}$ au degré près.
\end{questions}
\cote
\centering
\begin{tikzpicture}[otEspace,y={(0.62cm,0.5cm)},scale=2.1]
  \otPaveCoord{1}{1}{1}
  \coordinate (M) at (0,0.5,0); \coordinate (N) at (0.5,0,1);
  \fill[otGreen,opacity=0.18] (H) -- (M) -- (N) -- cycle;
  \otPaveAretes\otPaveNoms
  \draw[otGreenDark,line width=0.9pt] (H) -- (N);
  \draw[otGreenDark,line width=0.9pt,dash pattern=on 2.2pt off 1.6pt] (H) -- (M) -- (N);
  \node[otPtV] at (M) {}; \node[otSommet,anchor=north west] at (M) {$M$};
  \node[otPtV] at (N) {}; \node[otSommet,anchor=north] at (N) {$N$};
\end{tikzpicture}
\end{codecote}
\end{exercice}

\begin{exercice}[2]{Équations cartésiennes de plans}
L'espace est muni d'un repère orthonormé.
\begin{questions}
  \item Déterminer une équation du plan $\mathcal{P}_1$ passant par $A(2\pv0\pv1)$ et de vecteur normal $\vec n\,(1\pv-2\pv2)$.
  \item Le \textbf{plan médiateur} de $[AB]$ est l'ensemble des points équidistants de $A$ et de $B$. Justifier que c'est le plan passant
    par le milieu de $[AB]$ et de vecteur normal $\vect{AB}$, puis en donner une équation pour $A(1\pv2\pv3)$ et $B(3\pv0\pv-1)$.
  \item On considère le plan $\mathcal{P}:2x+y-z+3=0$. Donner un vecteur normal, vérifier que $M(1\pv-2\pv3)\in\mathcal{P}$, puis
    trouver deux vecteurs non colinéaires qui dirigent $\mathcal{P}$.
  \item Étudier la position relative de $\mathcal{P}$ et des plans $\mathcal{Q}:x-y+z=0$ et $\mathcal{R}:-4x-2y+2z+1=0$.
\end{questions}
\end{exercice}

\begin{exercice}[2]{Distance d'un point à un plan}
L'espace est muni d'un repère orthonormé. On considère le plan $\mathcal{P}:2x-y+2z-1=0$ et le point $M(3\pv1\pv4)$.
\begin{questions}
  \item Le point $M$ appartient-il à $\mathcal{P}$ ?
  \item Justifier que les points de la droite passant par $M$ et orthogonale à $\mathcal{P}$ s'écrivent $(3+2t\pv1-t\pv4+2t)$, avec $t$
    réel.
  \item En déduire les coordonnées du projeté orthogonal $H$ de $M$ sur $\mathcal{P}$, puis la distance de $M$ au plan $\mathcal{P}$.
  \item Vérifier que $Q(0\pv-1\pv0)$ appartient à $\mathcal{P}$ et calculer $MQ$. Est-ce cohérent ?
\end{questions}
\end{exercice}

\begin{exercice}[2]{Projeté orthogonal sur une droite}
L'espace est muni d'un repère orthonormé. On considère la droite $d$ passant par $B(0\pv1\pv1)$ et de vecteur directeur
$\vec u\,(1\pv2\pv2)$, ainsi que le point $M(3\pv1\pv4)$.
\begin{questions}
  \item Les points de $d$ s'écrivent $H_t=(t\pv1+2t\pv1+2t)$. Exprimer $\vect{MH_t}$ en fonction de $t$.
  \item Déterminer $t$ pour que $\vect{MH_t}\cdot\vec u=0$. En déduire le projeté orthogonal $H$ de $M$ sur $d$.
  \item Calculer la distance $MH$ du point $M$ à la droite $d$.
  \item Montrer que, pour tout réel $t$, $MH_t^2=9t^2-18t+18$. Retrouver ainsi que la distance $MH_t$ est minimale pour le point $H$.
\end{questions}
\end{exercice}

\begin{exercice}[2]{Deux plans perpendiculaires\enspace\otNouveau}
L'espace est muni d'un repère orthonormé. On considère les plans $\mathcal{P}_1:x-2y+2z-3=0$ et $\mathcal{P}_2:2x+2y+z-3=0$, ainsi que
les points $A(1\pv0\pv1)$ et $M(4\pv0\pv4)$.
\begin{questions}
  \item Montrer que les plans $\mathcal{P}_1$ et $\mathcal{P}_2$ sont perpendiculaires.
  \item Vérifier que $A$ appartient à $\mathcal{P}_1$ et à $\mathcal{P}_2$. Montrer que le vecteur $\vec u\,(-2\pv1\pv2)$ est orthogonal
    à un vecteur normal de chacun des deux plans. En déduire que la droite $\Delta$ passant par $A$ et dirigée par $\vec u$ est la
    droite d'intersection de $\mathcal{P}_1$ et de $\mathcal{P}_2$.
  \item Déterminer les coordonnées des projetés orthogonaux $H_1$ et $H_2$ du point $M$ sur $\mathcal{P}_1$ et sur $\mathcal{P}_2$,
    puis les distances $MH_1$ et $MH_2$.
  \item Montrer que le quadrilatère $MH_1AH_2$ est un carré.
  \item Vérifier que $\vect{AM}\cdot\vec u=0$. En déduire le projeté orthogonal de $M$ sur $\Delta$, puis la distance de $M$ à la
    droite $\Delta$.
\end{questions}
\end{exercice}

\rubrique{Problèmes}

\begin{exercice}[3]{Un angle variable dans un cube}
\begin{codecote}[0.6]
$ABCDEFGH$ est un cube d'arête $1$ ; on se place dans le repère orthonormé $\big(D\pv\vect{DA},\vect{DC},\vect{DH}\big)$. Pour
$x\in\intFF01$, le point $M$ de la diagonale $[DF]$ du cube est défini par $\vect{DM}=x\vect{DF}$, et l'on note $\theta(x)$ la mesure en
radians de l'angle $\widehat{EMB}$.
\begin{questions}
  \item Donner les coordonnées de $B$, $E$, $F$, puis montrer que $M(x\pv x\pv x)$.
  \item Montrer que $MB=ME$ et que $\vect{MB}\cdot\vect{ME}=3x^2-4x+1$. En déduire que
    $\cos\theta(x)=f(x)=\dfrac{3x^2-4x+1}{3x^2-4x+2}$.
\end{questions}
\cote
\centering
\begin{tikzpicture}[otEspace,scale=2.3]
  \otPaveCoord{1}{1}{1}
  \coordinate (M) at ($(D)!0.38!(F)$);
  \fill[otGreen,opacity=0.18] (E) -- (M) -- (B) -- cycle;
  \otPaveAretes\otPaveNoms
  \draw[otCache] (D) -- (F);
  \draw[otGreenDark,line width=0.9pt] (E) -- (B);
  \draw[otGreenDark,line width=0.9pt,dash pattern=on 2.2pt off 1.6pt] (E) -- (M) -- (B);
  \node[otPtV] at (M) {}; \node[otSommet,anchor=north] at (M) {$M$};
\end{tikzpicture}
\end{codecote}
\begin{questions}[start=3]
  \item Justifier que $3x^2-4x+2>0$ sur $\intFF01$, montrer que $f'(x)=\dfrac{2(3x-2)}{(3x^2-4x+2)^2}$, puis dresser le tableau de
    variations de $f$ en précisant $f(0)$, $f\left(\frac13\right)$, $f\left(\frac23\right)$ et $f(1)$.
  \item Pour quelles valeurs de $x$ l'angle $\widehat{EMB}$ est-il droit ? aigu ? obtus ? Pour quelle valeur de $x$ le triangle $MEB$
    est-il équilatéral ?
  \item Déterminer les valeurs minimale et maximale de l'angle $\widehat{EMB}$ lorsque $M$ parcourt $[DF]$.
\end{questions}
\end{exercice}

\begin{exercice}[3]{Un angle maximal dans un tétraèdre}
\begin{codecote}[0.62]
$OABC$ est un tétraèdre tel que $OA$, $OB$, $OC$ sont deux à deux perpendiculaires et de longueur $1$ : on travaille dans le repère
orthonormé $\big(O\pv\vect{OA},\vect{OB},\vect{OC}\big)$. Pour $x\in\intFF01$, le point $M$ de $[BC]$ est défini par
$\vect{BM}=x\vect{BC}$, et $\theta(x)$ est la mesure en radians de l'angle $\widehat{AMO}$.
\begin{questions}
  \item Justifier que $M(0\pv1-x\pv x)$, puis exprimer $\vect{MA}$ et $\vect{MO}$.
  \item Montrer que $\cos^2\theta(x)=g(x)=\dfrac{2x^2-2x+1}{2x^2-2x+2}$.
\end{questions}
\cote
\centering
\begin{tikzpicture}[scale=1.45]
  \coordinate (O) at (0,0); \coordinate (A) at (-0.85,-0.85); \coordinate (B) at (1.75,-0.15); \coordinate (C) at (0.1,1.6);
  \coordinate (M) at ($(B)!0.4!(C)$);
  \fill[otNavyTint,opacity=0.6] (A) -- (B) -- (C) -- cycle;
  \draw[otVu] (A) -- (B) -- (C) -- cycle;
  \draw[otCache] (O) -- (A) (O) -- (B) (O) -- (C);
  \draw[otGreenDark,line width=0.9pt] (A) -- (M);
  \draw[otGreenDark,line width=0.9pt,dash pattern=on 2.2pt off 1.6pt] (O) -- (M);
  \node[otPtV] at (M) {}; \node[otSommet,anchor=south west] at (M) {$M$};
  \node[otSommet,anchor=south east] at (O) {$O$}; \node[otSommet,anchor=north east] at (A) {$A$};
  \node[otSommet,anchor=west] at (B) {$B$}; \node[otSommet,anchor=south] at (C) {$C$};
\end{tikzpicture}
\end{codecote}
\begin{questions}[start=3]
  \item Calculer $\theta(0)$ et $\theta(1)$, et justifier que $0\leq\theta(x)\leq\frac{\pi}{2}$ pour tout $x\in\intFF01$.
  \item Étudier les variations de $g$ sur $\intFF01$. Pour quelle valeur de $x$ l'angle $\theta(x)$ est-il maximal ? Donner une valeur
    approchée de cet angle maximal, en degrés.
  \item Montrer que $\vect{MA}\cdot\vect{MB}=x(2x-1)$. Pour quelle position de $M$ distincte de $B$ le triangle $ABM$ est-il rectangle en $M$ ?
\end{questions}
\end{exercice}

\begin{exercice}[3]{La grande diagonale d'un cube}
\begin{codecote}[0.6]
$ABCDEFGH$ est un cube d'arête $1$ ; on travaille dans le repère orthonormé $\big(A\pv\vect{AB},\vect{AD},\vect{AE}\big)$.
\begin{questions}
  \item Montrer que le vecteur $\vect{AG}$ est normal au plan $(BDE)$, puis que ce plan a pour équation $x+y+z-1=0$.
  \item Déterminer les coordonnées du point $K$, intersection de la droite $(AG)$ et du plan $(BDE)$. Montrer que $K$ est le centre de
    gravité du triangle $BDE$.
  \item En déduire la distance du point $A$ au plan $(BDE)$.
  \item Calculer le volume du tétraèdre $ABDE$ de deux façons, et retrouver l'aire du triangle équilatéral $BDE$.
  \item Montrer que le plan $(CFH)$ est parallèle à $(BDE)$, et que les deux plans partagent la diagonale $[AG]$ en trois segments
    de même longueur.
\end{questions}
\cote
\centering
\begin{tikzpicture}[otEspace,scale=2.3]
  \otPaveCoord{1}{1}{1}
  \fill[otGreen,opacity=0.2] (B) -- (D) -- (E) -- cycle;
  \fill[otNavy,opacity=0.1] (C) -- (F) -- (H) -- cycle;
  \otPaveAretes\otPaveNoms
  \draw[otGreenDark,line width=0.8pt] (B) -- (E);
  \draw[otGreenDark,line width=0.8pt,dash pattern=on 2.2pt off 1.6pt] (B) -- (D) -- (E);
  \draw[otNavy,line width=0.7pt] (F) -- (H);
  \draw[otNavy,line width=0.7pt,dash pattern=on 2.2pt off 1.6pt] (F) -- (C) -- (H);
  \draw[otOrange,line width=0.8pt,dash pattern=on 2.2pt off 1.6pt] (A) -- (G);
  \coordinate (K) at (1/3,1/3,1/3);
  \node[otPtV] at (K) {}; \node[otSommet,anchor=north west] at (K) {$K$};
\end{tikzpicture}
\end{codecote}
\end{exercice}
